\section{Limitaciones}
\label{sec:limitations}

Enunciar las limitaciones con esta precisi\'on es en s\'i mismo parte
de la contribuci\'on argumentada en \S\ref{sec:related-work}: los
veredictos uno junto a otro de un agregador de reputaci\'on no llevan
un relato equivalente de d\'onde podr\'ian estar equivocados, lo que
los hace r\'apidos de consultar pero dif\'iciles de auditar. Cada
limitaci\'on siguiente se diagnostica, en cambio, a una causa concreta
y verificable en lugar de enunciarse de forma gen\'erica, de modo que
cada una ya est\'a mitigada de forma concreta o se acota como un
siguiente paso preciso (\S\ref{sec:future-work}) en lugar de quedar
como una advertencia abierta.

\begin{enumerate}

\item \label{lim:n3} \textbf{Conjunto de validaci\'on peque\~no (N=3).} Tres familias
  son pocas para una afirmaci\'on con potencia estad\'istica, y es el
  punto que m\'as merece escrutinio en este trabajo. El tama\~no no es
  solo una limitaci\'on de tiempo: los directores de esta tesis
  recomendaron expl\'icitamente acotarlo a tres familias para mantener
  el proyecto dentro de un alcance manejable como TFM, priorizando la
  profundidad de cada validaci\'on frente a dispersar el esfuerzo entre
  m\'as muestras. Tres mitigaciones compensan parcialmente ese tama\~no:
  los datos de referencia son ingenier\'ia inversa manual a nivel de
  funci\'on, verificada con Ghidra, para cada muestra, un list\'on
  considerablemente m\'as alto que el que tendr\'ia un conjunto m\'as
  grande pero superficial; una auditor\'ia completa de
  \code{static/scripts/} centrada en hardcoding espec\'ifico de muestra
  disfrazado de l\'ogica gen\'erica encontr\'o y elimin\'o cuatro
  instancias, cada una sustituida por un equivalente a nivel de clase y
  reverificada para no provocar regresi\'on en ninguna muestra
  validada; y el bucle de reenv\'io de RoningLoader
  (\S\ref{sec:resubmission}) punt\'ua de forma independiente 26
  componentes adicionales frente a los mismos datos de referencia,
  superficie de evaluaci\'on real adicional y no un truco estad\'istico.
  Ampliar a m\'as familias sigue siendo el siguiente paso de mayor valor
  (\S\ref{sec:future-work}).

\item \label{lim:attck-version} \textbf{La versi\'on de ATT\&CK fijada
  exige migraci\'on manual peri\'odica.} Todos los identificadores de
  t\'ecnica en esta pipeline, los datos de referencia y las capas de
  Navigator apuntan a MITRE ATT\&CK v19.2; las firmas de la comunidad
  de CAPE, en cambio, quedan fijadas a la versi\'on con la que se
  escribieron por \'ultima vez y se reasignan a su identidad actual
  solo en el momento de la curaci\'on (\S\ref{sec:dynamic}), no de
  forma autom\'atica en origen. Esta tesis migr\'o sus propios datos de
  referencia de v14 a v19 durante su desarrollo, un trabajo manual no
  trivial que a su vez sac\'o a la luz el fallo de calibraci\'on de
  confianza de \technique{T1070}/\technique{T1685.005} detallado en
  \S\ref{sec:fp-audit}. Cada futura revisi\'on de MITRE ATT\&CK exigir\'a
  repetir ese mismo esfuerzo, ya que nada en el dise\~no actual lo
  automatiza.

\item \label{lim:gt-table} \textbf{La extracci\'on de datos de
  referencia depende de la completitud de la tabla.} Un extractor basado en tablas markdown solo
  es tan completo como la propia tabla, incluso cuando la narrativa de
  un informe dice m\'as. \S\ref{sec:case-wsnake} es una instancia
  concreta, corregida en el documento fuente en lugar de sortearse en
  el extractor.

\item \label{lim:containment} \textbf{El sandboxing con contenci\'on
  prioritaria limita la observaci\'on de red/C2.} CAPE detona cada muestra contra la red
  simulada de INetSim, sin salida real, por dise\~no. Varias t\'ecnicas
  no detectadas en las tres muestras se agrupan en torno a
  comportamiento de exfiltraci\'on y proxy/t\'unel
  (\technique{T1041}, \technique{T1090}, \technique{T1571/T1572},
  \technique{T1046}) que una red simulada no puede provocar por
  completo. Esto es una propiedad estructural de un entorno de
  investigaci\'on con contenci\'on prioritaria, no un defecto a
  corregir.

\item \label{lim:key-ceiling} \textbf{La recuperaci\'on est\'atica de
  claves/configuraci\'on tiene un techo preciso.} El XOR de un solo byte es exhaustivamente
  forzable por fuerza bruta (128 claves), y as\'i se us\'o para
  recuperar la IP de C2 de RoningLoader sin conocimiento previo de la
  clave, pero solo tras a\~nadir dos filtros independientes de falsos
  positivos (\S\ref{sec:case-xor}). En cambio, la clave RC4 de una carga
  distinta, verificada de forma independiente mediante ingenier\'ia
  inversa manual con Ghidra, se ensambla en \emph{tiempo de
  ejecuci\'on} a partir de tres constantes combinadas mediante
  aritm\'etica de registros y nunca existe como secuencia de bytes
  contigua en el fichero: ning\'un m\'etodo est\'atico de patr\'on de
  bytes, por exhaustivo que sea, puede recuperar una clave que nunca se
  almacena como bytes. El l\'imite est\'a precisamente en si la clave
  existe como bytes en disco, no en una afirmaci\'on general de que
  ``el cifrado es dif\'icil''.

\item \label{lim:tool-constraints} \textbf{Limitaciones a nivel de
  herramienta.} La decodificaci\'on
  de cadenas basada en emulaci\'on de FLOSS agot\'o el tiempo en Akira,
  la muestra m\'as agresiva en anti-an\'alisis; esto se registr\'o
  expl\'icitamente en lugar de faltar en silencio. La recuperaci\'on
  AES/ChaCha20 solo tiene \'exito frente a un nonce/IV cero. Cualquier
  nueva capacidad de recuperaci\'on por fuerza bruta ahora se valida,
  por pol\'itica expl\'icita, frente a una segunda muestra antes de
  confiar en ella, tras detectarse as\'i la cascada de falsos positivos
  de la clave XOR (\S\ref{sec:case-xor}) en lugar de al primer \'exito.

\item \label{lim:reproducibility} \textbf{Reproducibilidad.} La \S\ref{sec:deployment} describe el
  despliegue completo con \code{docker-compose}; el \'unico hueco que
  merece repetirse aqu\'i es que la m\'aquina virtual Windows del
  perfil \code{sandbox}
  necesita una configuraci\'on manual \'unica del host (libvirt/KVM),
  documentada con sus propias notas de soluci\'on de problemas, ya que
  no puede crearse solo con \code{docker-compose}.

\item \label{lim:rule-coverage} \textbf{La cobertura de reglas
  est\'aticas es deliberadamente estrecha.} Aproximadamente 85 de las 474 t\'ecnicas relevantes para
  Windows en MITRE tienen una regla est\'atica dedicada, con la
  justificaci\'on del alcance detallada en \S\ref{sec:coverage}.

\item \label{lim:network-attribution} \textbf{La captura de red
  din\'amica no tiene atribuci\'on por proceso.} La lista \code{network.hosts} de CAPE refleja el
  tr\'afico de toda la m\'aquina virtual del sandbox, no solo el del
  propio proceso de la muestra, por lo que la telemetr\'ia de fondo del
  sistema operativo hu\'esped puede aparecer indistinguible de las
  conexiones reales de la muestra (\S\ref{sec:case-ioc-noise}). La
  mitigaci\'on aplicada es una lista de permitidos estrecha con las IP
  espec\'ificas confirmadas individualmente que realmente contaminaban
  las detonaciones de estas tres muestras, no un filtro general de ASN
  o de reputaci\'on; una muestra reci\'en detonada podr\'ia mostrar la
  misma clase de contaminaci\'on bajo una IP distinta que la lista
  actual todav\'ia no cubre.

\end{enumerate}
